\section{Parameterized algorithms and kernels}\label{sec:algorithms}

The hardness map of Section~\ref{sec:complexity} is dense, yet it leaves
positive windows, and this section fills them. We give an \XP\ algorithm
in the number of opened facilities, the branch-and-bound
refinement we implement, the fixed-parameter algorithm that pins down the
size-parameter dichotomy, the routing certificates that bridge the
enumeration to the exact formulation, and the two sides of the
kernelization question. A short remark on width parameters closes the
account, and Table~\ref{tab:landscape} then collects the whole landscape.

\subsection{An \XP\ algorithm in the number of opened
facilities}\label{sec:xp}

Once a design is fixed the cost is determined by the routing oracle of
Proposition~\ref{prop:oracle}, so the only combinatorial freedom left is
which facilities to open. Bounding that freedom by $k$ leaves at most
$e\,n^{k}$ candidate supports, with $e$ the base of the natural
logarithm, and each is priced by one oracle call.
The next theorem promotes the brute-force membership of
Observation~\ref{obs:xp} to a running time whose exponent is linear in
$k$, and the branch-and-bound that follows stays within the same
$n^{k+O(1)}$ regime while pruning aggressively in practice.

The \emph{support} of a design is the set of facilities it opens. We
identify each set $S\subseteq I\cup J$ with the design that opens exactly
the facilities of $S$, and under this identification we write $v(S)$ for
its routing value, $f(S)=\sum_{i\in S\cap I}f_i$ and
$g(S)=\sum_{j\in S\cap J}g_j$ for its two fixed charges, and
$\mathrm{cost}(S)=f(S)+g(S)+v(S)$ for its total cost. On an infeasible
support the convention $v(S)=+\infty$ of Section~\ref{sec:prelim}
gives $\mathrm{cost}(S)=+\infty$.

\begin{theorem}\label{thm:xp}
Write $T_{\mathrm{orac}}$ for the time of one routing-oracle call and
normalize $k\le n$. Then $\mathrm{MP\text{-}TSCFLP}(B,k)$ is decidable,
and the least cost $\OPT_k$ over feasible designs opening at most $k$
facilities is computable, in time $O\bigl(n^{k}\cdot
(T_{\mathrm{orac}}+n|L|)+|K||L|\bigr)$ for $k\ge 1$, and in time
$O(T_{\mathrm{orac}}+n|L|+|K||L|)$
for $k=0$, the additive term reading the demand totals once. In particular the problem lies in \XP\ in the parameter $k$.
\end{theorem}
\begin{proof}
A feasible design opening at most $k$ facilities has a
support
$S\subseteq I\cup J$ with $|S|\le k$, its cost is the fixed charge of $S$
plus the routing value $v(S)$, and by Proposition~\ref{prop:oracle} the
routing value is attained, so each feasible support realizes exactly its
cost and $\OPT_k$ is the minimum of $\mathrm{cost}(S)$ over feasible
supports of size at most $k$. When no such support is feasible we put
$\OPT_k=+\infty$, and the correct answer is then \textsc{no}. The totals $D_1,\dots,D_{|L|}$ are computed once in $O(|K||L|)$ time
and reused by every feasibility test, which is the additive term of the
statement. Enumerate
every such $S$, test feasibility
by Proposition~\ref{prop:feas} in $O(n|L|)$ time against the
precomputed totals, and price the feasible
ones by one oracle call. The number of supports is
$\sum_{t=0}^{k}\binom{n}{t}\le\sum_{t=0}^{k}\frac{n^{t}}{t!}
\le n^{k}\sum_{t\ge 0}\frac{1}{t!}=e\,n^{k}$
for $k\le n$ and $k\ge 1$. The middle inequality uses $n^{t}\le n^{k}$
for $t\le k$, valid since $n\ge 1$, and extending the sum adds
nonnegative terms. Each support
costs one $O(n|L|)$ feasibility test plus at most one oracle call, and
$T_{\mathrm{orac}}$ is left as a stated quantity rather than bounded
below, since an implementation may share preprocessed data across calls,
giving $O\bigl(n^{k}(T_{\mathrm{orac}}+n|L|)+|K||L|\bigr)$ in total. For
$k=0$ only the empty design is priced. That design is feasible exactly
when every $D_l=0$, since for it the left-hand sums of (F1) and (F2) are
empty, so one feasibility test and at most one oracle call suffice,
within $O(T_{\mathrm{orac}}+n|L|+|K||L|)$ counting the totals. Comparing the minimum
against $B$ decides $\mathrm{MP\text{-}TSCFLP}(B,k)$.
\end{proof}

\subsection{A practical branch-and-bound refinement}\label{sec:bb}

Theorem~\ref{thm:xp} settles the classification, and the rest is
implementation. The enumeration it counts is wasteful whenever a
partial design already cannot be
completed within the budget, and the branch-and-bound of
Algorithm~\ref{alg:xpbb}, the code the computational study runs,
closes such subtrees early. Its principal device
is a covering count. Because both stages are complete, feasibility of a
product needs only enough open capacity on each side, and a sorted prefix
determines the fewest facilities that supply a demand from a nonnegative
capacity vector.

\begin{lemma}[covering count]\label{lem:cover}
Let $a_1,\dots,a_m\ge 0$ and $D\ge 0$, and let
$a_{(1)}\ge\cdots\ge a_{(m)}$ be the decreasing order. Some $S$ with
$|S|\le s$ has $\sum_{i\in S}a_i\ge D$ if and only if
$\sum_{r\le s}a_{(r)}\ge D$. When $\sum_{r\le m}a_{(r)}<D$ no such $s$
exists and we set $k^\ast(a;D)=+\infty$. Otherwise the least such $s$,
written $k^\ast(a;D)$, is computable by one sort and one prefix scan in
$O(m\log m)$ time, the same scan detecting the infinite case.
\end{lemma}
\begin{proof}
Both directions follow from majorization by the sorted prefix. If
$\sum_{r\le s}a_{(r)}\ge D$ the set of the $s$ largest entries witnesses
the forward direction. Conversely, order any $S$ decreasingly. Its $r$th
element is at most $a_{(r)}$, because the $r$ largest elements of $S$ are
$r$ distinct entries of $a$, each at least that element, so $a$ has at
least $r$ entries at least as large. Summing gives
$\sum_{i\in S}a_i\le\sum_{r\le|S|}a_{(r)}$.
Hence if some $S$ with $|S|\le s$ has $\sum_{i\in S}a_i\ge D$ then
$D\le\sum_{i\in S}a_i\le\sum_{r\le|S|}a_{(r)}\le\sum_{r\le s}a_{(r)}$, the
last step by $a\ge 0$. The prefix sums are nondecreasing in $s$, so the
least $s$ with $\sum_{r\le s}a_{(r)}\ge D$ is found by one sort and one
scan.
\end{proof}

Algorithm~\ref{alg:xpbb} searches over partial designs, and we name the
quantities its nodes carry before giving the listing. A node is a pair
$(O,C)$ in which $O$ is the set of facilities forced open and $C$ the set
forced closed, so that $F=(I\cup J)\setminus(O\cup C)$ collects the free
facilities and $r=k-|O|$ is the remaining cardinality. For a product $l$
let
\[
s_I(l;O,C)=k^\ast\!\Bigl((b_{il})_{i\in F\cap I};\,\max\bigl(0,\,D_l-\textstyle
\sum_{i\in O\cap I}b_{il}\bigr)\Bigr),
\]
the covering count of Lemma~\ref{lem:cover} over the free factory
capacities against the demand left after crediting the already-open
factories, and let
\[
s_J(l;O,C)=k^\ast\!\Bigl((p_{jl})_{j\in F\cap J};\,\max\bigl(0,\,D_l-\textstyle
\sum_{j\in O\cap J}p_{jl}\bigr)\Bigr)
\]
be the analogous count for (F2) over the free depot capacities. The
quantities of line 2 are the root specializations
$k^\ast_I(l)=k^\ast\bigl((b_{il})_{i\in I};D_l\bigr)=s_I(l;\emptyset,\emptyset)$
and
$k^\ast_J(l)=k^\ast\bigl((p_{jl})_{j\in J};D_l\bigr)=s_J(l;\emptyset,\emptyset)$.
Let $v^{\uparrow}=v(O\cup F)$ be the routing value of the all-open
completion, with the convention $v=+\infty$ of
Section~\ref{sec:prelim}, and let $LB(O,C)=f(O)+g(O)+v^{\uparrow}$.
Finally, a facility of $O$ is \emph{unused} by a flow when the flow
carries zero on every arc incident to that facility in every product's
network. Three prunes act on these quantities, and after the listing we
show each is sound and then that the search returns $\OPT_k$.
Figure~\ref{fig:bbtree} places the three prunes on a single branching
step.

% Requires: \usetikzlibrary{positioning,arrows.meta}
\begin{figure}[H]
  \centering
  \begin{tikzpicture}[
      >={Stealth[length=1.8mm]},
      res/.style={draw=black!70,line width=0.4pt,fill=black!6,
        rounded corners=2pt,align=center,inner sep=3.5pt,
        font=\footnotesize},
      note/.style={font=\scriptsize,text=black!55,align=center},
      branch/.style={draw=black,line width=0.5pt,->},
      lead/.style={draw=black!55,line width=0.35pt,dashed}
    ]
    \node[res] (P) at (0,0)
      {$(O,C)$\\
       \textcolor{black!60}{\scriptsize free set $F$,\quad $r=k-|O|$}};
    \node[res] (L) at (-2.5,-2.1)
      {$(O\cup\{e\},C)$\\
       \textcolor{black!60}{\scriptsize open $e$}};
    \node[res] (R) at (2.5,-2.1)
      {$(O,C\cup\{e\})$\\
       \textcolor{black!60}{\scriptsize close $e$}};
    \draw[branch] (P) -- (L);
    \draw[branch] (P) -- (R);
    % P1 acts on the popped node, covering counts against r
    \node[note,anchor=east] (p1) at (-3.1,0.5)
      {P1, covering\\ $\max_l s_I+\max_l s_J>r$};
    \draw[lead] (p1.east) -- (P.west);
    % P2 acts on the popped node, admissible bound
    \node[note,anchor=west] (p2) at (3.1,0.5)
      {P2, bound\\ $LB(O,C)\ge\min(\mathrm{best},B+1)$};
    \draw[lead] (p2.west) -- (P.east);
    % P3 acts only at a leaf
    \node[note] (p3) at (0,-3.0)
      {P3, dominance, only at a leaf, where $F=\varnothing$ and the
       oracle flow leaves some facility of $O$ unused};
  \end{tikzpicture}
  \caption{One branching step of Algorithm~\ref{alg:xpbb}. A node
  carries the forced-open set $O$, the forced-closed set $C$, the free
  set $F$ and the remaining cardinality $r$, and branching on a free
  facility $e$ produces the two children along the solid arrows. The
  dashed leaders attach P1 and P2 to the popped node they test, while
  P3 fires only at a leaf, so the three prunes act exactly on the named
  quantities (Lemma~\ref{lem:cover},
  Proposition~\ref{prop:bbcorrect}).}
  \label{fig:bbtree}
\end{figure}


\begin{algorithm}[H]
\caption{Branch-and-bound for $\mathrm{MP\text{-}TSCFLP}(B,k)$
(Theorem~\ref{thm:xp}).}
\label{alg:xpbb}
\begin{pseudolines}
\item $k\gets\min(k,n)$ \algcomment{normalization, no design opens more}
\item \kw{for each} product $l$ sort $(b_{il})_i$ and $(p_{jl})_j$ and
  compute $k^\ast_I(l),k^\ast_J(l)$ \algcomment{Lemma~\ref{lem:cover}}
\item \kw{if}
  $\max_l k^\ast_I(l)+\max_l k^\ast_J(l)>k$ \kw{then} \kw{return}
  \textsc{no} with $\OPT_k=+\infty$
  \algcomment{a priori bound $k^\ast$}
\item $\mathrm{best}\gets+\infty$; push the root $(O,C)=(\emptyset,\emptyset)$
\item \kw{while} the stack is nonempty \kw{do}
\item \algin pop $(O,C)$; let $F$ be the free facilities and
  $r\gets k-|O|$; \kw{if} $r<0$ \kw{then} continue
\item \algin \kw{if}
  $\max_l s_I(l;O,C)+\max_l s_J(l;O,C)>r$ \kw{then} continue
  \algcomment{prune P1, Lemma~\ref{lem:cover}}
\item \algin evaluate the all-open completion by the oracle to get
  $v^{\uparrow}$ and $LB(O,C)$
\item \algin \kw{if} $v^{\uparrow}=+\infty$ \kw{or}
  $LB(O,C)\ge\min(\mathrm{best},B+1)$ \kw{then} continue
  \algcomment{prune P2}
\item \algin \kw{if} $F=\emptyset$ \kw{then}
  \algcomment{leaf, the design is $S=O$}
\item \algin\algin \kw{if} the oracle solution returned at the leaf
  leaves some facility of $O$ unused \kw{then} continue
  \algcomment{prune P3}
\item \algin\algin $\mathrm{best}\gets\min(\mathrm{best},LB(O,C))$
\item \algin \kw{else} pick $e\in F$ of least index and push
  $(O\cup\{e\},C)$ and $(O,C\cup\{e\})$
\item \kw{if} $\mathrm{best}<+\infty$ \kw{and} $\mathrm{best}\le B$
  \kw{then} \kw{return} \textsc{yes} \kw{else} \kw{return} \textsc{no},
  either way reporting $\mathrm{best}$
  \algcomment{run with $B=\infty$ to read $\OPT_k$ off either branch}
\end{pseudolines}
\end{algorithm}

\emph{P1, covering.} Every completion of $(O,C)$ opens $O$ together with
some free facilities. For such a completion $S$ and a product $l$, (F1)
reads $\sum_{i\in S\cap I}b_{il}\ge D_l$, which states that the free
factories of $S$ supply at least the residual demand
$\max(0,D_l-\sum_{i\in O\cap I}b_{il})$, so Lemma~\ref{lem:cover},
applied to $(b_{il})_{i\in F\cap I}$ against that residual demand, shows
that $S$ adds at least $s_I(l;O,C)$ free factories. One open set must
satisfy (F1) for all
products at once, so it adds at least $\max_l s_I(l;O,C)$ free factories,
and likewise at least $\max_l s_J(l;O,C)$ free depots. Factories and depots
are disjoint, so the two increments add, and when their sum exceeds $r$ no
completion has cardinality at most $k$, so P1 discards $(O,C)$ soundly.
An infinite count falls under the same reading, since
$s_I(l;O,C)=+\infty$ means that even opening every free factory leaves
(F1) violated for $l$, so no completion is feasible and the prune fires
for every finite $r$.

\emph{P2, bound.} Any completion opens a set $S$ with $O\subseteq
S\subseteq O\cup F$. Its fixed charge is at least $f(O)+g(O)$ by
nonnegativity of the charges, and its routing
value is at least $v(O\cup F)=v^{\uparrow}$ because $S\le O\cup F$
componentwise and $v$ is monotone (Proposition~\ref{prop:mono}). Hence
$LB(O,C)$ is a lower bound on the cost of \emph{every} completion, even one
opening more than $k$ facilities, which is all the bound needs. P2 discards
$(O,C)$ when $LB(O,C)\ge\min(\mathrm{best},B+1)$, and either threshold
justifies the discard. If $LB(O,C)\ge\mathrm{best}$ then no completion
improves the incumbent. If $LB(O,C)\ge B+1$ then no completion meets the
budget, because the data are nonnegative integers and the routing values
are integers by Proposition~\ref{prop:oracle}, so every cost is an
integer and a cost at most $B$ is a cost strictly below $B+1$. When
$v^{\uparrow}=+\infty$ the same monotonicity gives
$v(S)\ge v(O\cup F)=+\infty$ for every completion $S$, so every
completion is infeasible under the convention of
Section~\ref{sec:prelim} and discarding is again sound. At a
leaf, where $F=\emptyset$, one has $LB(O,C)=\mathrm{cost}(O)$.

\emph{P3, dominance.} When the oracle solution returned at the leaf
leaves an open facility of $O$ unused, closing that facility gives a
design of strictly smaller cardinality. The oracle solution returned at
the leaf carries nothing on any arc of the
closed facility, so it remains a feasible routing of the smaller design,
whence the routing value does not increase, and the fixed charge drops
by the nonnegative charge of that facility. The smaller design is
therefore feasible at no greater cost, so $O$ is dominated and P3
discards it.

With the three prunes justified, the remaining question is whether the
surviving search still reaches an optimum.

\begin{proposition}[correctness]\label{prop:bbcorrect}
Algorithm~\ref{alg:xpbb} decides $\mathrm{MP\text{-}TSCFLP}(B,k)$, and when
run with $B=+\infty$ it returns $\mathrm{best}=\OPT_k$.
\end{proposition}
The full argument appears in Appendix~\ref{app:bb}. Its three steps justify
the early return of line 3 through Lemma~\ref{lem:cover}, bound
$\mathrm{best}$ below by $\OPT_k$ because only feasible designs within
the cardinality budget ever reach a leaf, and bound it above by tracking
a minimum-cardinality optimum down the tree, checking that no prune ever
discards the node carrying it.

The prunes only discard nodes, and the tree they act on is itself
small. Branching always selects the free facility of least index, so at
every node the decided set $O\cup C$ is a prefix of the facility order
and the node is determined by the prefix length $m$ together with
$O\subseteq\{1,\dots,m\}$. Children are pushed only from nodes with
$|O|\le k$, so every node ever created has $|O|\le k+1$ and the tree
holds at most
$\sum_{m=0}^{n}\sum_{t=0}^{k+1}\binom{m}{t}=O(n^{k+2})$ nodes, since
$\sum_{t=0}^{k+1}\binom{m}{t}\le\sum_{t=0}^{k+1}n^{t}\le 2n^{k+1}$
for $n\ge 2$ by the geometric series and the cases $n\le 1$ are
trivial. Each
node performs the sorts and prefix scans of P1 and at most one oracle
call, so the whole search runs in
$O\bigl(n^{k+2}(T_{\mathrm{orac}}+n|L|\log n)+|K||L|\bigr)$ time, with
the additive term again reading the demand totals once, a
polynomial factor above the enumeration bound of Theorem~\ref{thm:xp}
and within the same $n^{k+O(1)}$ regime, which is all the
classification uses. On the seeded validation suite
recorded in the accompanying code, the reference implementation visited
roughly ninety per cent fewer nodes than the plain enumeration and returned
the brute-force optimum on every instance. The
linear exponent is essentially the best one can hope for, since under the
Exponential Time Hypothesis no algorithm decides the problem in time
$f(k)\cdot N^{o(k)}$ (Theorem~\ref{thm:eth}), so the enumeration of
supports is tight up to the constant in the exponent.

One relation among the prunes has consequences that
Section~\ref{sec:exp-q1} meets again on the machine, so we record it
here.

\begin{observation}[the covering prune in plain optimization]
\label{obs:p1redundant}
Run with $k=n$, prune P1 discards no node beyond the infeasibility
branch of P2, since at any node $r=n-|O|=|F|+|C|\ge|F|$. If the
all-open
completion $O\cup F$ is feasible, then for every product the free
facilities of each side supply that side's residual demand, so
Lemma~\ref{lem:cover} gives $\max_l s_I(l;O,C)\le|F\cap I|$ and
$\max_l s_J(l;O,C)\le|F\cap J|$, the two maxima add to at most
$|F|\le r$, and P1 does not fire. When P1 fires, therefore, the
all-open completion is infeasible, so $v^{\uparrow}=+\infty$ and the
infeasibility branch of P2 discards the same node at the price of one
oracle call. Under a binding cardinality $k<n$ the subsumption fails,
since a node may admit a feasible all-open completion while no
completion within the remaining budget $r$ exists, a situation P1 can detect
and the infeasibility branch of P2 cannot.
\end{observation}

\subsection{Fixed-parameter tractability and the size-parameter
dichotomy}\label{sec:fpt}

The parameter $k$ counts openings, whereas the four dimensions
$|I|,|J|,|K|,|L|$ measure the instance itself, and only the two facility
sides carry tractability. Enumerating designs rather than supports gives a
fixed-parameter algorithm in $|I|+|J|$, and coupling it with the hardness
slices of Section~\ref{sec:complexity} pins the exact border that
Figure~\ref{fig:dichotomy} maps over the sixteen subsets.

% Requires: \usetikzlibrary{positioning,fit}
\begin{figure}[H]
  \centering
  \resizebox{0.98\textwidth}{!}{%
  \begin{tikzpicture}[
      x=1cm,y=1.35cm,
      every node/.style={font=\scriptsize},
      nd/.style={draw=black!70,line width=0.4pt,fill=black!4,rounded corners=1.5pt,
        minimum width=8mm,minimum height=5mm,inner sep=2.5pt},
      fpt/.style={draw,line width=0.8pt,fill=black!18,rounded corners=1.5pt,
        font=\scriptsize\bfseries,minimum width=8mm,minimum height=5mm,inner sep=2.5pt},
      weak/.style={draw=black!70,line width=0.5pt,dashed,fill=black!4,rounded corners=1.5pt,
        minimum width=8mm,minimum height=5mm,inner sep=2.5pt},
      ed/.style={draw=black!40,line width=0.3pt}
    ]
    % rank 0
    \node[nd] (e) at (4.5,0) {$\varnothing$};
    % rank 1
    \node[nd] (I) at (1.5,1) {I}; \node[nd] (J) at (3.5,1) {J};
    \node[nd] (K) at (5.5,1) {K}; \node[nd] (L) at (7.5,1) {L};
    % rank 2
    \node[fpt] (IJ) at (0.5,2) {IJ}; \node[nd] (IK) at (2.1,2) {IK};
    \node[nd] (IL) at (3.7,2) {IL}; \node[nd] (JK) at (5.3,2) {JK};
    \node[weak] (JL) at (6.9,2) {JL$^\dagger$}; \node[nd] (KL) at (8.5,2) {KL};
    % rank 3
    \node[fpt] (IJK) at (1.5,3) {IJK}; \node[fpt] (IJL) at (3.5,3) {IJL};
    \node[weak] (IKL) at (5.5,3) {IKL$^\dagger$}; \node[weak] (JKL) at (7.5,3) {JKL$^\dagger$};
    % rank 4
    \node[fpt] (IJKL) at (4.5,4) {IJKL};
    % cover edges (drawn behind)
    \begin{scope}[on background layer]
      \foreach \a/\b in {e/I,e/J,e/K,e/L,
        I/IJ,I/IK,I/IL, J/IJ,J/JK,J/JL, K/IK,K/JK,K/KL, L/IL,L/JL,L/KL,
        IJ/IJK,IJ/IJL, IK/IJK,IK/IKL, IL/IJL,IL/IKL,
        JK/IJK,JK/JKL, JL/IJL,JL/JKL, KL/IKL,KL/JKL,
        IJK/IJKL,IJL/IJKL,IKL/IJKL,JKL/IJKL}{\draw[ed] (\a) -- (\b);}
    \end{scope}
    % rank labels
    \foreach \r/\t in {0/{$|P|=0$},1/{$1$},2/{$2$},3/{$3$},4/{$4$}}{
      \node[anchor=east,font=\scriptsize\itshape,text=black!55] at (-0.6,\r) {\t};}
    % legend
    \begin{scope}[shift={(0,-1.15)}]
      \node[fpt,minimum width=4mm,minimum height=3.5mm] (lf) at (0,0) {};
      \node[anchor=west] at (0.35,0) {FPT ($\{$I,J$\}\subseteq P$)};
      \node[nd,minimum width=4mm,minimum height=3.5mm] (ls) at (3.9,0) {};
      \node[anchor=west] at (4.25,0) {para-NP-hard (strong)};
      \node[weak,minimum width=4mm,minimum height=3.5mm] (lw) at (0,-0.75) {};
      \node[anchor=west] at (0.35,-0.75) {$\dagger$ para-NP-hard (weak under binary encoding)};
    \end{scope}
  \end{tikzpicture}%
  }
  \caption{The size-parameter dichotomy of Theorem~\ref{thm:dichotomy}
  over the sixteen subsets $P\subseteq\{|I|,|J|,|K|,|L|\}$, arranged by
  cardinality as a Hasse diagram whose edges are the cover relations.
  The problem is fixed-parameter tractable exactly on the four
  upward-closed sets that contain both facility layers I and J, shaded
  and shown in bold (Theorem~\ref{thm:fpt}), and para-NP-hard on the
  remaining twelve. Hardness is strong except at the three sets marked
  with a dagger, whose witnessing slices are only weakly NP-hard under a
  binary encoding (Theorem~\ref{thm:partition},
  Corollary~\ref{cor:partitionmirror}).}
  \label{fig:dichotomy}
\end{figure}


\begin{theorem}[fixed-parameter tractability]\label{thm:fpt}
$\mathrm{MP\text{-}TSCFLP}(B)$, its optimization form and its cardinality
version are solvable in time $2^{|I|+|J|}\cdot\mathrm{poly}(N)$, where $N$
is the input size. Hence the problem is \FPT\ in the combined parameter
$(|I|,|J|)$, and a fortiori under any parameter set containing both.
\end{theorem}
\begin{proof}
Enumerate all $2^{|I|+|J|}$ designs, test each by
Proposition~\ref{prop:feas}, price the feasible ones by
Proposition~\ref{prop:oracle}, filter by $\sum y+\sum z\le k$ when the
cardinality bound is present, and return the least cost, compared
against $B$ for the decision versions, reporting infeasibility through
the value $+\infty$ when no design passes the test. Correctness is immediate from the
attainment in Proposition~\ref{prop:oracle}, each step per design runs
in time polynomial in $N$, and the function $2^{|I|+|J|}$ reads only the
two facility counts, which is the definition of \FPT\ in the combined
parameter. The a fortiori clause follows because $2^{|I|+|J|}$ is also a
computable function of any parameter set that contains both facility
counts.
\end{proof}

The synthesis theorem states the border cleanly. For a parameter set $P$
drawn from the four dimensions the problem is fixed-parameter tractable
exactly when $P$ contains both facility sides, and otherwise it is hard
already at a fixed slice.

\begin{theorem}[dichotomy]\label{thm:dichotomy}
Let $P\subseteq\{|I|,|J|,|K|,|L|\}$ parameterize
$\mathrm{MP\text{-}TSCFLP}(B)$. If $\{|I|,|J|\}\subseteq P$ the problem is
\FPT, unconditionally. If $\{|I|,|J|\}\not\subseteq P$ it is
para-\NP-hard, unconditionally, because some assignment of fixed values to
the parameters of $P$ gives an \NP-hard slice, and then no \XP\ algorithm
exists unless $\Ppoly=\NP$. Assuming $\Ppoly\ne\NP$, therefore, the
problem is \FPT\ if and only if $\{|I|,|J|\}\subseteq P$. The same
statement transfers verbatim to $\mathrm{MP\text{-}TSCFLP}(B,k)$ by
appending $k:=n$.
\end{theorem}
\begin{proof}
The positive half is Theorem~\ref{thm:fpt}. For the negative half it
suffices to exhibit, for each of the twelve subsets $P$ with
$\{|I|,|J|\}\not\subseteq P$, an \NP-hard slice that fixes every parameter
of $P$ to the constant one, since a problem \NP-hard at a fixed parameter
value is para-\NP-hard and an \XP\ algorithm would decide that slice in
polynomial time. Six slices from Section~\ref{sec:complexity} cover all
twelve. The slice of Theorem~\ref{thm:cover} fixes $|I|=|L|=1$ and
witnesses $\emptyset$, $\{|I|\}$, $\{|L|\}$ and $\{|I|,|L|\}$. The slice of
Theorem~\ref{thm:products} fixes $|J|=|K|=1$ and witnesses $\{|J|\}$,
$\{|K|\}$ and $\{|J|,|K|\}$, while its mirror, Corollary~\ref{cor:mirror},
fixes $|I|=|K|=1$ and witnesses $\{|I|,|K|\}$. The slice of
Theorem~\ref{thm:numeric} fixes $|K|=|L|=1$ and witnesses $\{|K|,|L|\}$.
The weak slice of Theorem~\ref{thm:partition} fixes $|J|=|K|=|L|=1$ and
witnesses $\{|J|,|L|\}$ and the triple $\{|J|,|K|,|L|\}$, and its depot
mirror, Corollary~\ref{cor:partitionmirror}, fixes $|I|=|K|=|L|=1$ and
witnesses the remaining triple $\{|I|,|K|,|L|\}$. Each of the twelve
subsets with $\{|I|,|J|\}\not\subseteq P$ is contained in one of the six
fixed sets $\{|I|,|L|\}$, $\{|J|,|K|\}$, $\{|I|,|K|\}$, $\{|K|,|L|\}$,
$\{|J|,|K|,|L|\}$, $\{|I|,|K|,|L|\}$, and fixing a superset of $P$ to
constants fixes $P$, so the case analysis is complete. Appending $k:=n$
preserves the answer, because no design opens more than $n$ facilities,
so the appended bound is vacuous and the two problems coincide instance
by instance, and it leaves the fixed values of $P$ untouched.
\end{proof}

Two qualifications attach to the negative half. Para-\NP-hardness
asks only for \NP-hardness of a slice under the standard binary encoding,
and weak hardness already is \NP-hardness, so the weakly hard slices count
towards the dichotomy exactly as the strongly hard ones do. The encoding
matters only when we ask whether the border survives unary numbers. There
the positive half is untouched, and the negative half survives for the
nine parameter sets whose witnessing slices are strongly \NP-hard, while
for $\{|J|,|L|\}$ and the two triples $\{|I|,|K|,|L|\}$ and
$\{|J|,|K|,|L|\}$ the witnessing slices turn polynomial under unary data,
through the pseudo-polynomial programs of Theorem~\ref{thm:partition}
and Corollary~\ref{cor:partitionmirror},
and no other hard slice is known, so their unary status stays open.

\subsection{Certificates and a bridge to the exact
formulation}\label{sec:certificates}

The dichotomy settles what the parameters buy, and this subsection
turns the same oracle into certificates a mixed-integer solver can
consume. Each oracle call solves one transportation problem per product, and its
linear-programming dual both certifies the routing value and supplies a
cut that a Benders master can reuse. The fact that carries the
subsection is that the dual polyhedron depends only on the data, never
on the design, so a single dual point
yields an inequality valid for every design whose block is feasible.
Figure~\ref{fig:bendersflow} traces the pipeline the proposition
formalizes.

% Requires: \usetikzlibrary{positioning,arrows.meta}
\begin{figure}[H]
  \centering
  \resizebox{0.97\textwidth}{!}{%
  \begin{tikzpicture}[
      >={Stealth[length=1.8mm]},
      res/.style={draw=black!70,line width=0.4pt,fill=black!6,
        rounded corners=2pt,align=center,inner sep=3.5pt,
        font=\footnotesize},
      flow/.style={draw=black,line width=0.5pt,->}
    ]
    \node[res,text width=2.1cm] (D) at (0,0)
      {design $(y,z)$\\
       \textcolor{black!60}{\scriptsize from master or enumeration}};
    \node[res,text width=2.6cm] (O) at (3.6,0)
      {routing oracle\\
       \textcolor{black!60}{\scriptsize one transportation block per
       product $l$}};
    \node[res,text width=3.0cm] (C) at (7.5,0)
      {primal and dual certificates\\
       \textcolor{black!60}{\scriptsize optimal flow and
       $u^\ast\in\Delta_l$, or a ray}};
    \node[res,text width=2.6cm] (K) at (11.2,0)
      {disaggregated cut\\
       \textcolor{black!60}{\scriptsize optimality on $\theta_l$, or
       feasibility}};
    \draw[flow] (D) -- (O);
    \draw[flow] (O) -- (C);
    \draw[flow] (C) -- (K);
  \end{tikzpicture}%
  }
  \caption{From a design to a cut, the pipeline of
  Proposition~\ref{prop:duals}. Each oracle call of
  Proposition~\ref{prop:oracle} prices the design one product at a
  time, returns an optimal flow together with a dual point
  $u^\ast\in\Delta_l$ certifying the block value, and that dual point
  yields the disaggregated Benders optimality cut, valid for every
  design with a feasible block because $\Delta_l$ never depends on $(y,z)$, while an
  infeasible block yields a recession ray and a feasibility cut
  instead.}
  \label{fig:bendersflow}
\end{figure}


\begin{proposition}[routing certificate and Benders cuts]\label{prop:duals}
Fix a product $l$ and let $\Delta_l$ be the dual polyhedron of its
transportation block, which is independent of $(y,z)$. If the block is
feasible then some $u^\ast\in\Delta_l$ attains the dual value
$\mu_l(y,z)$, and the pair of an optimal flow and $u^\ast$ certifies
$\mu_l(y,z)$ in time $O(|I||J|+|J||K|)$. Further, writing $\theta_l$ for
a variable through which a master problem represents the block value, for
every $u=(\alpha,\beta,\sigma,\pi)\in\Delta_l$ the inequality
\[
\theta_l\ \ge\ \sum_k q_{kl}\,\alpha_k
      -\sum_i b_{il}\,\sigma_i\,y_i-\sum_j p_{jl}\,\pi_j\,z_j
\]
is valid for every design whose block $l$ is feasible, and it holds with
equality at the design that generated $u^\ast$. Two explicit recession
rays of $\Delta_l$ certify infeasibility, and their objectives are
positive exactly when (F1), respectively (F2), of
Proposition~\ref{prop:feas} fails.
\end{proposition}
The proof, given in Appendix~\ref{app:duals}, writes the block dual explicitly.
No dual constraint mentions the design, which yields the independence of
$\Delta_l$, weak duality yields the displayed inequality, strong duality
supplies the attaining point $u^\ast$, and the two recession rays are
exhibited coordinate by coordinate and matched to the failures of (F1)
and (F2).

These are precisely the per-product optimality cuts of a disaggregated
Benders decomposition, and the recession rays are its feasibility cuts.
The certificates the enumeration consumes are the objects the Benders
decomposition emits, so each oracle call in Algorithm~\ref{alg:xpbb} can
return, beside its value, the very cut that a branch-and-Benders-cut
master would add at that design. The proposition certifies the routing
value and the validity of the cuts, not the global optimum, which still
needs the master's own bounds.

\subsection{Kernelization}\label{sec:kernel}

Preprocessing is the last algorithmic question, and it divides cleanly.
On the positive side, customers with identical cost columns aggregate
exactly, a bound on costs and total demands yields a polynomial kernel
outright, and a bound on costs and budget still yields a polynomial
compression through the weight-reduction technique of
\citet{FrankTardos1987}. On the negative side, products do not
aggregate, and no polynomial kernel exists in
the winning parameter unless the polynomial hierarchy collapses. The first
reduction is safe and cheap.

\begin{proposition}[customer aggregation]\label{prop:aggregation}
If customers $k\ne k'$ have identical cost columns, that is
$d_{jkl}=d_{jk'l}$ for all $j,l$, then merging $k'$ into $k$ and adding
their demands preserves feasibility, routing value and cost of every
design, hence the optimal value, the set of optimal designs and the
answers of both decision versions. Feasible complete solutions
map in both directions under the explicit merge and split maps of the
proof, which never raise the cost and preserve the optimal routing
value of every design, and consequently one
may assume $|K|\le\tau$, where $\tau$
is the number of distinct second-stage cost columns, computable in
polynomial time.
\end{proposition}
We prove this in Appendix~\ref{app:agg} by exhibiting the two maps directly. The
merge adds the flow rows of the two customers and preserves every
constraint and every cost term. The split trims any surplus of the
merged row and expands it by the northwest-corner procedure, which
preserves feasibility at no cost increase, so the optimal routing
values agree design by design, and a companion reduction bounds the
capacities without touching any design.

\begin{remark}[aggregation up to additive shifts]\label{rem:addshift}
The exact-equality hypothesis can be relaxed at the price of a budget
shift. Suppose the columns of $k$ and $k'$ differ by a product-wise
constant, $d_{jk'l}=d_{jkl}+\alpha_l$ for all $j,l$. In a routing
trimmed as in the proof of Proposition~\ref{prop:aggregation},
customer $k'$ receives exactly $q_{k'l}$ units of product $l$, each
paying $\alpha_l$ more than under the column of $k$ at the same
depot. Replacing the column of $k'$ by that of $k$ therefore shifts
the
routing value of every feasible design by exactly
$-\sum_l\alpha_l q_{k'l}$, and the shift reaches the optima because
they are attained at trimmed routings, trimming never raising the
cost in either instance. The merge of
Proposition~\ref{prop:aggregation} then applies to the now identical
columns once the budget is replaced by
$B-\sum_l\alpha_l q_{k'l}$, the answer being \textsc{no} outright
when that quantity is negative, since costs are nonnegative. We keep
$\tau$ as defined, the refined count of columns distinct up to
product-wise shifts not being needed elsewhere.
\end{remark}

\begin{observation}[capacity capping]\label{obs:capping}
Replacing every capacity $b_{il}$ by $\min(b_{il},D_l)$ and every
$p_{jl}$ by $\min(p_{jl},D_l)$ preserves feasibility, routing value and
cost of every design, hence the optimum and both decision answers.
\end{observation}
Appendix~\ref{app:cap} proves this by trimming any routing until every
product moves exactly $D_l$ across each stage, at no cost increase,
after which no capacity beyond $D_l$ can carry flow and the cap changes
nothing.

The symmetric move on products is unsound. In analogy with
Proposition~\ref{prop:aggregation}, merging two products with identical
cost columns would add their demands and add their capacity columns
entrywise while keeping the shared cost columns, and that operation lets
the instance transfer capacity between the merged products. Take two
factories with charges $f=(0,10)$, one free depot of ample
capacity, one customer, all transport costs zero, and two products of unit
demand with capacities $b_{\cdot l}=(2,0)$ and $b_{\cdot l'}=(0,2)$.
In the original instance (F1) for $l'$ forces the second factory open at
charge $10$, while every other charge and every transport cost vanishes, so
the optimum $10$ is both forced and attained. The merged product has
demand $2$ against capacities $(2,2)$ and the
first factory alone serves it at cost $0$, so the merge can lower
the value. A variant shows the merge can also repair infeasibility. In
the same setting give the two products the capacity columns
$b_{\cdot l}=(2,0)$ and $b_{\cdot l'}=(0,0)$, so that the original
instance is infeasible, (F1) failing for $l'$ at every design, while the
merged product has demand $2$ against the column $(2,0)$ and the first
factory serves it. The natural rule that adds demands and capacity
columns is therefore unsound, and the paper advances no product-side
aggregation. A one-sided bound on the data still buys genuine
preprocessing, a kernel on one side and a compression on the other,
along the two routes that
Figure~\ref{fig:ftpipeline} lays out.

% Requires: \usetikzlibrary{positioning,arrows.meta}
\begin{figure}[H]
  \centering
  \resizebox{0.99\textwidth}{!}{%
  \begin{tikzpicture}[
      >={Stealth[length=1.8mm]},
      res/.style={draw=black!70,line width=0.4pt,fill=black!6,
        rounded corners=2pt,align=center,inner sep=3.5pt,
        font=\footnotesize},
      flow/.style={draw=black,line width=0.5pt,->}
    ]
    \node[res,text width=1.7cm] (I) at (0,0)
      {instance\\
       \textcolor{black!60}{\scriptsize costs, totals
       $\le\kappa^{c^\star}$}};
    \node[res,text width=2.2cm] (A) at (2.95,0)
      {customer aggregation\\
       \textcolor{black!60}{\scriptsize Prop.~\ref{prop:aggregation},
       $|K|\le\tau$}};
    \node[res,text width=2.0cm] (C) at (5.95,0)
      {capacity capping\\
       \textcolor{black!60}{\scriptsize Obs.~\ref{obs:capping}, caps
       $\le\max_l D_l$}};
    \node[res,text width=2.3cm] (S) at (8.95,0)
      {bounded coefficient forms\\
       \textcolor{black!60}{\scriptsize basic flows,
       $|a|\le\Delta_{\max}$}};
    \node[res,text width=2.0cm] (F) at (11.95,0)
      {Frank--Tardos\\
       \textcolor{black!60}{\scriptsize $\Phi\mapsto\Phi'$,
       $\mathrm{poly}(\kappa)$ bits}};
    \node[res,text width=2.0cm] (L) at (14.85,0)
      {sign-condition language\\
       \textcolor{black!60}{\scriptsize equivalent output}};
    \draw[flow] (I) -- (A);
    \draw[flow] (A) -- (C);
    \draw[flow] (C) -- (S);
    \draw[flow] (S) -- (F);
    \draw[flow] (F) -- (L);
  \end{tikzpicture}%
  }
  \caption{The compression pipeline of
  Proposition~\ref{prop:ftcompress}. Aggregation and capping
  (Prop.~\ref{prop:aggregation}, Obs.~\ref{obs:capping}) bound the
  entities and the numeric entries by $\mathrm{poly}(\kappa)$, total
  unimodularity of the oracle networks of
  Proposition~\ref{prop:oracle} turns acceptance into sign conditions
  on integer forms with coefficients of magnitude at most
  $\Delta_{\max}$, and the Frank--Tardos step replaces the cost vector
  $\Phi$ by an equivalent $\Phi'$ of $\mathrm{poly}(\kappa)$ bits, so
  the bounded subclass compresses in $\kappa$.}
  \label{fig:ftpipeline}
\end{figure}


\begin{proposition}[kernel and compression under one-sided
bounds]\label{prop:ftcompress}
Fix $c^\star$ and let $\kappa=|I|+|J|+|L|+\tau$. The subclass of
$\mathrm{MP\text{-}TSCFLP}(B)$ whose opening and transport costs
$f,g,c,d$ and whose total demands $D_1,\dots,D_{|L|}$ are integers at
most $\kappa^{c^\star}$, with capacities and budget unrestricted,
admits a polynomial kernel in $\kappa$, an equivalent instance of
$\mathrm{MP\text{-}TSCFLP}(B)$ itself of $\mathrm{poly}(\kappa)$
bits. The subclass with $\max(f,g,c,d,B)\le\kappa^{c^\star}$ and
capacities and demands unrestricted admits a
polynomial compression in $\kappa$ into a fixed sign-condition
language, which is not itself a location problem. The first subclass is nontrivial.
The strongly hard slice with $|K|=|L|=1$ of Theorem~\ref{thm:numeric}
has costs in $\{0,1\}$ and a single customer whose demand, the only
product total, is $|I|\,(|J|+1)\le\kappa^{2}$, so the slice lies in the
subclass for every $c^\star\ge 2$.
\end{proposition}
The proof, given in Appendix~\ref{app:ft}, is elementary for the first
subclass. Aggregation and capping leave every datum except the budget
within $\kappa^{c^\star}$, the all-open design decides feasibility
outright, and its cost, itself polynomially bounded in $\kappa$, caps
the budget, so the output is an ordinary instance of the problem. No
such cap exists in the second subclass, whose large numbers are the
capacities and demands, and there acceptance becomes a Boolean
combination of sign conditions on integer linear forms with
coefficients that total unimodularity bounds by a fixed power of
$\kappa$, which the
Frank--Tardos theorem compresses to
$\mathrm{poly}(\kappa)$ bits preserving every sign at once.

The matching negative result completes the account. Two cross-compositions
of \textsc{Set Cover} show that, unless
$\NP\subseteq\mathrm{coNP}/\mathrm{poly}$, the parameter of the dichotomy
admits no polynomial kernel, and
the same construction reaches the augmented parameters and the unary
encoding.

\begin{theorem}[no polynomial kernel]\label{thm:nokernel}
Unless $\NP\subseteq\mathrm{coNP}/\mathrm{poly}$,
$\mathrm{MP\text{-}TSCFLP}(B)$ parameterized by $|I|+|J|$ admits no
polynomial compression, and in particular no polynomial kernel, already on
the subclass $|I|=|L|=1$ and even with unary data. The exclusion transfers
to the parameters $|I|+|J|+|L|$ and $|I|+|J|+|K|$, to the cardinality
version $\mathrm{MP\text{-}TSCFLP}(B,k)$, and to the subclass $|K|=1$.
\end{theorem}
\begin{proof}
We give an OR-cross-composition of \textsc{Set Cover} \citep{Karp1972}
into $\mathrm{MP\text{-}TSCFLP}(B)$ parameterized by $|I|+|J|$, which by
the framework of \citet{BodlaenderJansenKratsch2014} rules out a
polynomial compression unless $\NP\subseteq\mathrm{coNP}/\mathrm{poly}$.
The framework asks for a polynomial equivalence relation, and we take
the relation $R$ that declares two well-formed instances equivalent when
they agree on the universe size $|U|$, the family size $m$ and the
target $t$, with $|U|\ge 1$ and $1\le t\le m$. Four further classes are
reserved, one for the malformed inputs, one for the well-formed
instances with empty universe, one for those with nonempty universe and
$t=0$, and one for those with nonempty universe and $t>m$.
Each instance falling in one of these four classes is decided
individually in polynomial time, a malformed input being rejected by
the syntactic check that recognizes it, an empty universe being covered
by the empty subfamily, a target of zero over a nonempty universe
admitting no cover, and a target above $m$ admitting one exactly when
the whole family covers the universe. On inputs from those classes the composition
therefore answers the OR of those individual answers, emitting a fixed
yes- or no-instance realizing it. Deciding $R$ compares three numbers, and
the triple $(|U|,m,t)$ takes at most cubically many values in the total
encoding length of the input list, so
$R$ is a polynomial equivalence relation in the sense of the framework.

Take now $T$ \textsc{Set Cover} instances $X_1,\dots,X_T$ from one
nontrivial class, sharing $|U|\ge 1$, $m$ and $t$ with $1\le t\le m$,
and assume
$T=2^{s}$ with $s=\lceil\log T\rceil$ by duplicating instances,
duplicating once more when $T=1$ so that $s\ge 1$. Identify the copies
with the strings of $\{0,1\}^{s}$ and write $c_a$ for the $a$th bit of
the copy index $c$. The budget is $B=s(t+1)+t$, the amplifier is
$Q=B+1$, and the customers comprise $T|U|$ copy customers and $s$ guard
customers, one guard per index bit, each of demand $Q$, so a customer
forced onto arcs of unit cost pays at least $Q>B$ even under divisible
service, exactly the accounting of Theorem~\ref{thm:cover}. The
composed instance has a single product and a single factory of opening
cost zero, capacity $(T|U|+s)\,Q$ equal to the total demand, and zero
first-stage costs, so the whole construction lives in the depot layer
and the second-stage matrix, and $|I|=|L|=1$ throughout.

\emph{Depots.} Use $m$ set-depots $D_1,\dots,D_m$ of opening cost $1$,
shared across all copies, and $s$ pairs of selector depots
$(P_a^{0},P_a^{1})$, one pair per index bit, each of opening cost $t+1$.
Every depot has capacity equal to the total demand. Sharing the
set-depots is sound, and no instance is merged, because the
copy-specific set membership is carried not by the depots but by the
cost matrix to the copy-specific customers introduced next. Opening
$P_a^{b}$ reads bit $a$ of the selected index as $b$. A guard customer
per pair reaches either selector of its own pair at cost $0$ and every
other depot at cost $1$. A pair with neither selector open therefore
leaves its guard on unit-cost arcs at a cost of at least $Q$, so at
least one selector per pair must open. The opening charges cap the
number of open selectors, since $s+1$ open selectors already cost
$(s+1)(t+1)=s(t+1)+t+1=B+1>B$, so
jointly with the guards a within-budget design opens exactly one
selector per pair, fixing a single index $\iota\in\{0,1\}^{s}$. Because
all routing costs are nonnegative, the selector charges leave for the
set-depot openings a residual budget of exactly $B-s(t+1)=t$.

\emph{Customers.} For every copy $c$ and element $u\in U$ create a
customer $(c,u)$ with demand $Q$. The cost from set-depot $D_j$ to
$(c,u)$ is $0$ when $u$ belongs to the $j$th set of copy $c$ and $1$
otherwise, so the whole of copy $c$'s instance lives in these costs. The
cost from selector $P_a^{b}$ to $(c,u)$ is $0$ when $c_a\ne b$, and $1$
otherwise.

\emph{Correctness.} Fix the selected index $\iota$. A copy $c\ne\iota$
differs from $\iota$ in some bit $a$, so the open selector
$P_a^{\iota_a}$ serves all of copy $c$ for free, and the non-selected
copies cost nothing beyond the selectors already open. The copy $\iota$
agrees with the pattern on every bit, so every selector arc into its
customers costs one, and each $(\iota,u)$, carrying demand $Q$, must be
served through a set-depot whose set in copy $\iota$ contains $u$, since
any other service costs at least $Q>B$. Each open set-depot costs $1$
against the residual budget $t$, so at most $t$ set-depots are open.
Every element $u$ therefore has an open set-depot whose set in copy
$\iota$ contains it, and those sets form a cover of
$X_\iota$ of size at most $t$. Conversely, if some $X_\iota$ has a cover
of size at most $t$, open the factory, the selector pattern of $\iota$ and
that cover. Every depot has capacity equal to the total demand, so with
the free factory and these depots open (F1) and (F2) hold and the design
is feasible by Proposition~\ref{prop:feas}. Completeness of the two
stages then lets every customer draw its demand through a zero-cost
depot. The guards use the open selector of their own pair, the customers
of a copy $c\ne\iota$ use the open selector $P_a^{\iota_a}$ at a bit $a$
where $c$ disagrees with $\iota$, as in the forward direction, and the
customers of copy $\iota$ use the depots of the cover, so the routing
cost is zero and the total cost is the fixed charge
$s(t+1)+|C|\le s(t+1)+t=B$. A within-budget design therefore exists if and
only if the OR of the $T$ instances is a yes-instance.

\emph{Parameter and robustness.} The depots number $m+2s$ and the
factories one, so $|I|+|J|=O(m+\log T)$, polynomially bounded in the
largest input and in $\log T$ exactly as the framework requires. The
construction writes $O((T|U|+s)(m+2s))$ cost entries and numbers of
$O(\log(T|U|m))$ bits, so it runs in time polynomial in
$\sum_c\lvert X_c\rvert$, which is the framework's time bound. The
largest numbers are the capacities, equal to the total demand
$(T|U|+s)\,Q$. Here $Q=B+1=s(t+1)+t+1\le(s+1)(m+1)=O(m\log T)$ since
$t\le m$, and $s\le T\le T|U|$ since $|U|\ge 1$, so the total demand is
$O(T\,|U|\,m\log T)$, hence polynomial in the output size,
so the composition is robust to a unary encoding, and the single product
places it in the subclass $|I|=|L|=1$ while inflating $|K|$.

\emph{Transfer to $|K|=1$.} The dual construction keeps one depot and one
customer and moves the composition to the factory and product sides,
mirroring the primal exactly as Corollary~\ref{cor:mirror} mirrors
Theorem~\ref{thm:products}. Introduce a product $(c,u)$ of demand $Q$ for every
copy $c$ and element $u$. Introduce $m$ shared set-factories
$F_1,\dots,F_m$ of opening cost $1$, whose first-stage cost to product
$(c,u)$ is $0$ when $u$ lies in the $i$th set of copy $c$ and $1$
otherwise, together with $s$ pairs of selector factories of opening cost $t+1$,
whose cost to $(c,u)$ is $0$ when $c_a\ne b$ for the pair's bit $a$ and
read value $b$, and $1$ otherwise. Add a guard product per pair, of
demand $Q$, free through either selector factory of its pair and costing
$1$ elsewhere. The budget is again $B=s(t+1)+t$ with amplifier $Q=B+1$.
The single depot and the single customer are free, the
customer demands $Q$ units of every product, and all second-stage costs
vanish. Every factory and the depot hold, in each product separately,
capacity equal to the total demand, so
feasibility reads through (F1) and (F2) as before, and since the second
stage carries zero cost the entire penalty sits on the first stage. The
paragraphs on depots and correctness of the primal now apply verbatim under
the substitution of set-factory for set-depot, selector factory for
selector depot, product $(c,u)$ for customer $(c,u)$, guard product for
guard customer, and first-stage cost for second-stage cost, with (F1),
(F2) and the budget arithmetic unchanged. In the transferred
construction the single depot receives every product through the first
stage and forwards it to the single customer, so constraint \eqref{eq:C2}
reads, for each product, first-stage inflow at least second-stage
outflow, and the routing that serves each product through its designated
factory satisfies it with equality. The orientation of the balance
constraint is therefore respected and the budget arithmetic is
unchanged. A within-budget design
therefore selects one index $\iota$ and covers copy $\iota$'s products
with at most $t$ set-factories, and conversely a cover of size at most
$t$ yields a design of cost at most $B$. This
keeps $|K|=1$ while inflating $|L|$, with $|I|+|J|=(m+2s)+1=O(m+\log T)$
and all numbers polynomial in the output, so it excludes a polynomial
kernel on the subclass $|K|=1$ too.

Finally, the augmented parameters
match the two compositions as follows. Adding $|L|$ leaves
$|I|+|J|+|L|$ polynomially bounded in the primal composition, where
$|L|=1$, and adding $|K|$ leaves $|I|+|J|+|K|$ polynomially bounded in
the dual composition, where $|K|=1$, so each augmented parameter
inherits the exclusion from its composition. Appending the cardinality
bound $k=1+s+t$ is likewise redundant. The total demand is positive, so
(F1) forces the single factory of the primal open and (F2) forces the
single depot of the dual open, and every within-budget design opens that
forced facility, exactly one selector per pair and at most $t$
set-facilities, hence at most $1+s+t=k$ facilities in all. The witness
of the converse direction opens $1+s+|C|\le 1+s+t=k$ of them, so the
bound thins the set of witnesses and breaks
neither direction.
\end{proof}

Fixed-parameter tractability in $|I|+|J|$ still yields an exponential
kernel, by the standard equivalence between fixed-parameter
tractability and kernelization \citep[e.g.][]{Cygan2015}, since on
instances larger than a threshold depending on the parameter the \FPT\
algorithm runs in polynomial time and a trivial equivalent instance can
be output, while below the threshold the instance is its own kernel. What
Theorem~\ref{thm:nokernel} rules out is only the polynomial one, and
whether a polynomial kernel exists for the full parameter
$|I|+|J|+|K|+|L|$ with general numbers, where the yes-certificate
comparison becomes genuinely bilinear, is left open.

\subsection{Width parameters}\label{sec:width}

A reader versed in algorithmic meta-theorems might hope to route the
tractable side through graph width, and this last remark closes that door
for the representation we analyse.
The natural incidence graph carries no useful width, and over that
representation the obstruction
is arithmetic rather than structural.

\begin{observation}[width is uninformative]\label{obs:cliquewidth}
We first fix the graph the claim is about.
For every instance the layered incidence graph on $I\cup J\cup K$ is the
complete bipartite graph $K_{J,\,I\cup K}$, since
$K_{I,J}\cup K_{J,K}=K_{J,\,I\cup K}$. Its clique-width is at most $2$,
witnessed by the expression that introduces every vertex of $J$ with
label $1$, introduces every vertex of $I\cup K$ with label $2$, and
performs a single join between the two labels.

Its treewidth
\citep{RobertsonSeymour1986} equals
$\min(|J|,|I|+|K|)$ once both sides are nonempty, and is therefore
unbounded over the class. Indeed, for $K_{m',n'}$ with
$1\le m'\le n'$, the bags formed by the smaller side together with one
vertex of the larger side, one bag per such vertex and arranged along a
path, give a decomposition of width $m'$. The matching lower bound
follows by contracting a matching of size $m'-1$ between the two sides.
The contracted
vertices become pairwise adjacent, and the remaining vertex of the
smaller side and any remaining vertex of the larger side are adjacent to
all of them and to each other, so $K_{m',n'}$ has a $K_{m'+1}$ minor,
and since $K_{m'+1}$ has treewidth $m'$ and treewidth never increases
under minors \citep{RobertsonSeymour1986}, the treewidth is exactly
$m'=\min(m',n')$.

The logical consequence closes the observation.
The treewidth-one star slice with $|J|=1$ is already strongly
\NP-hard, by Theorem~\ref{thm:products} (which fixes $|J|=|K|=1$). Were
the problem expressible in the
$\mathrm{MSO}_1$ optimization logic over the unweighted incidence graph
used here, the clique-width meta-theorem
\citep{CourcelleMakowskyRotics2000} together with the bound of at most
two would force a polynomial algorithm on the whole class,
contradicting strong \NP-hardness unless $\Ppoly=\NP$, so that route is
closed. The capacity and demand constraints compare sums of numbers,
which that logic over that graph cannot express, and over the
representation used here the obstruction is therefore arithmetic rather
than structural. This persists even when every capacity, demand and
cost lies in $\{0,1,m+1\}$ and the budget is at most $m$, where $m$ is
the size of the set family in the construction of
Theorem~\ref{thm:products}. Richer relational structures, weighted or
counting logics, or representations carrying products and constraints
as objects are not addressed by this observation, and their width
status simply remains open.
\end{observation}
